\documentclass[10pt,a4paper]{amsart}
\usepackage{amsmath,amssymb,amsthm,graphicx}
\usepackage{verbatim}
\newenvironment{reference}{\comment}{\endcomment}
%\newenvironment{reference}{}{}
\newenvironment{slogan}{\comment}{\endcomment}
\newenvironment{history}{\comment}{\endcomment}

% For commutative diagrams we use Xy-pic
\usepackage[all]{xy}

% We use 2cell for 2-commutative diagrams.
\xyoption{2cell}
% Theorem environments.
%
\theoremstyle{plain}
\newtheorem{theorem}[subsection]{Theorem}
\newtheorem{proposition}[subsection]{Proposition}
\newtheorem{lemma}[subsection]{Lemma}

\theoremstyle{definition}
\newtheorem{definition}[subsection]{Definition}
\newtheorem{example}[subsection]{Example}
\newtheorem{exercise}[subsection]{Exercise}
\newtheorem{situation}[subsection]{Situation}

\theoremstyle{remark}
\newtheorem{remark}[subsection]{Remark}
\newtheorem{remarks}[subsection]{Remarks}

\numberwithin{equation}{subsection}

% Macros
%
\def\lim{\mathop{\mathrm{lim}}\nolimits}
\def\colim{\mathop{\mathrm{colim}}\nolimits}
\def\Spec{\mathop{\mathrm{Spec}}}
\def\Hom{\mathop{\mathrm{Hom}}\nolimits}
\def\Ext{\mathop{\mathrm{Ext}}\nolimits}
\def\SheafHom{\mathop{\mathcal{H}\!\mathit{om}}\nolimits}
\def\SheafExt{\mathop{\mathcal{E}\!\mathit{xt}}\nolimits}
\def\Sch{\mathit{Sch}}
\def\Mor{\mathop{\mathrm{Mor}}\nolimits}
\def\Ob{\mathop{\mathrm{Ob}}\nolimits}
\def\Sh{\mathop{\mathit{Sh}}\nolimits}
\def\NL{\mathop{N\!L}\nolimits}
\def\CH{\mathop{\mathrm{CH}}\nolimits}
\def\proetale{{pro\text{-}\acute{e}tale}}
\def\etale{{\acute{e}tale}}
\def\QCoh{\mathit{QCoh}}
\def\Ker{\mathop{\mathrm{Ker}}}
\def\Im{\mathop{\mathrm{Im}}}
\def\Coker{\mathop{\mathrm{Coker}}}
\def\Coim{\mathop{\mathrm{Coim}}}

% Boxtimes
%
\DeclareMathSymbol{\boxtimes}{\mathbin}{AMSa}{"02}

%
% Macros for moduli stacks/spaces
%
\def\QCohstack{\mathcal{QC}\!\mathit{oh}}
\def\Cohstack{\mathcal{C}\!\mathit{oh}}
\def\Spacesstack{\mathcal{S}\!\mathit{paces}}
\def\Quotfunctor{\mathrm{Quot}}
\def\Hilbfunctor{\mathrm{Hilb}}
\def\Curvesstack{\mathcal{C}\!\mathit{urves}}
\def\Polarizedstack{\mathcal{P}\!\mathit{olarized}}
\def\Complexesstack{\mathcal{C}\!\mathit{omplexes}}
% \Pic is the operator that assigns to X its picard group, usage \Pic(X)
% \Picardstack_{X/B} denotes the Picard stack of X over B
% \Picardfunctor_{X/B} denotes the Picard functor of X over B
\def\Pic{\mathop{\mathrm{Pic}}\nolimits}
\def\Picardstack{\mathcal{P}\!\mathit{ic}}
\def\Picardfunctor{\mathrm{Pic}}
\def\Deformationcategory{\mathcal{D}\!\mathit{ef}}
\usepackage{cleveref}
\crefname{originalreference}{original reference}{original references}
\Crefname{originalreference}{Original reference}{Original references}
\makeatletter
\newlabel{lemma-compute-what-it-should-be}{{0CTD}{source}{}{sourcepointer-1}{}}
\newlabel{lemma-compute-what-it-should-be@cref}{{[originalreference][0][]0CTD}{[1][0][]source}{}{}{}}
\makeatother
\crefname{originalreference}{original reference}{original references}
\Crefname{originalreference}{Original reference}{Original references}
\makeatletter
\newlabel{lemma-finitely-presented}{{0CTH}{source}{}{sourcepointer-2}{}}
\newlabel{lemma-finitely-presented@cref}{{[originalreference][0][]0CTH}{[1][0][]source}{}{}{}}
\makeatother
\crefname{originallemma}{original lemma}{original lemmas}
\Crefname{originallemma}{Original lemma}{Original lemmas}
\makeatletter
\newlabel{more-algebra-lemma-limit-henselian}{{0CT7}{source}{}{sourcepointer-3}{}}
\newlabel{more-algebra-lemma-limit-henselian@cref}{{[originallemma][0][]0CT7}{[1][0][]source}{}{}{}}
\makeatother
\crefname{originalreference}{original reference}{original references}
\Crefname{originalreference}{Original reference}{Original references}
\makeatletter
\newlabel{situation-existence}{{08B7}{source}{}{sourcepointer-4}{}}
\newlabel{situation-existence@cref}{{[originalreference][0][]08B7}{[1][0][]source}{}{}{}}
\makeatother
\crefname{originallemma}{original lemma}{original lemmas}
\Crefname{originallemma}{Original lemma}{Original lemmas}
\makeatletter
\newlabel{spaces-more-morphisms-lemma-split-off-proper-part-henselian}{{0CWI}{source}{}{sourcepointer-5}{}}
\newlabel{spaces-more-morphisms-lemma-split-off-proper-part-henselian@cref}{{[originallemma][0][]0CWI}{[1][0][]source}{}{}{}}
\makeatother
\title{Shrinking the open so the support of the degree zero part is proper over the base}
\author{The Stacks Project Authors}
\date{stacks-project:spaces-flat.tex:local-snapshot-20261012; GNU FDL 1.2 or later}
\begin{document}
\maketitle
\noindent\textit{Editorial note.} All delivered mathematical text is the author's own, copied from the pinned Stacks Project chapter file spaces-flat.tex; no statement or proof step was rewritten, reordered or removed.\par\medskip
\noindent\textit{Editorial note.} The counted proof is the author's complete argument: henselian level pairs, support identification after base change, the henselian open and closed split-off of the support, and the shrinking of W by the part disjoint from the special fibre\par\medskip
\noindent\textit{Editorial note.} The label lemma-proper-support names the counted statement, and the proof environment immediately following it in the source is delivered verbatim as its proof.\par\medskip
\noindent\textit{Editorial note.} The counted statement is carried by the author's own lemma environment in the chapter source.\par\medskip
\noindent\textit{Allowed standard prerequisites (editorial).}\begin{itemize}
\item Henselianity of the pair (A, I\_n) for I\_n the kernel of A -> A\_n
\item The scheme theoretic support of a finitely presented module as a closed subspace
\item The stage-wise identification of the support after base change to A\_n from the finitely presented lemma of the situation
\item The henselian split-off lemma writing a closed subspace as a disjoint union of an open and closed part proper over the base and a part disjoint from the closed fibre
\end{itemize}\medskip
\ifdefined\hypertarget\hypertarget{sourcepointer-1}{}\fi
\noindent\textit{Source reference.} Original reference 0CTD: sources/stacks-spaces-flat/pinned-original-source.tex, source lines 2388--2388. Prerequisite pointer invoked at proof line(s) 2575; it names a standard fact of the same chapter that the counted argument uses, and every step of the counted proof built on it is retained verbatim in the delivered body.\par\medskip
\ifdefined\hypertarget\hypertarget{sourcepointer-2}{}\fi
\noindent\textit{Source reference.} Original reference 0CTH: sources/stacks-spaces-flat/pinned-original-source.tex, source lines 2508--2508. Prerequisite pointer invoked at proof line(s) 2576; it names a standard fact of the same chapter that the counted argument uses, and every step of the counted proof built on it is retained verbatim in the delivered body.\par\medskip
\ifdefined\hypertarget\hypertarget{sourcepointer-3}{}\fi
\noindent\textit{Source reference.} Original lemma 0CT7: sources/stacks-chow/tags.txt, source lines 16057--16057. Prerequisite pointer invoked at proof line(s) 2583; it names a standard fact of a sibling chapter that the counted argument uses, and every step of the counted proof built on it is retained verbatim in the delivered body.\par\medskip
\ifdefined\hypertarget\hypertarget{sourcepointer-4}{}\fi
\noindent\textit{Source reference.} Original reference 08B7: sources/stacks-spaces-flat/pinned-original-source.tex, source lines 2362--2362. Prerequisite pointer invoked at proof line(s) 2574; it names a standard fact of the same chapter that the counted argument uses, and every step of the counted proof built on it is retained verbatim in the delivered body.\par\medskip
\ifdefined\hypertarget\hypertarget{sourcepointer-5}{}\fi
\noindent\textit{Source reference.} Original lemma 0CWI: sources/stacks-chow/tags.txt, source lines 16173--16173. Prerequisite pointer invoked at proof line(s) 2592; it names a standard fact of a sibling chapter that the counted argument uses, and every step of the counted proof built on it is retained verbatim in the delivered body.\par\medskip
\begin{lemma}
\label{lemma-proper-support}
In Situation \ref{situation-existence} let $K$ be as in
Lemma \ref{lemma-compute-what-it-should-be}. Let $W \subset X$
be as in Lemma \ref{lemma-finitely-presented}.
Set $\mathcal{F} = H^0(K)|_W$. Then, after possibly shrinking the open $W$,
the support of $\mathcal{F}$ is proper over $A$.
\end{lemma}

\begin{proof}
Fix $n \geq 1$. Let $I_n = \Ker(A \to A_n)$.
By More on Algebra, Lemma \ref{more-algebra-lemma-limit-henselian}
the pair $(A, I_n)$ is henselian.
Let $Z \subset W$ be the scheme theoretic support of $\mathcal{F}$.
This is a closed subspace as $\mathcal{F}$ is of finite presentation.
By part (3) of Lemma \ref{lemma-finitely-presented}
we see that $Z \times_{\Spec(A)} \Spec(A_n)$
is equal to the support of $\mathcal{F}_n$ and hence
proper over $\Spec(A/I)$.
By More on Morphisms of Spaces, Lemma
\ref{spaces-more-morphisms-lemma-split-off-proper-part-henselian}
we can write $Z = Z_1 \amalg Z_2$ with $Z_1, Z_2$ open and
closed in $Z$, with $Z_1$ proper
over $A$, and with $Z_1 \times_{\Spec(A)} \Spec(A/I_n)$
equal to the support of $\mathcal{F}_n$.
In other words, $|Z_2|$ does not meet $X_0$.
Hence after replacing $W$ by $W \setminus Z_2$ we obtain the lemma.
\end{proof}

\end{document}
